\documentclass[11pt]{article}
\usepackage[margin=1in]{geometry}
\usepackage{amsmath}
\usepackage{booktabs}

\title{Do Kernel Machines Need Feature Scaling? Evidence from Breast Cancer Diagnosis}
\author{Paper2Proof canary set\thanks{A small paper written for the Paper2Proof evaluation set (eval set A).}}
\date{2026}

\begin{document}
\maketitle

\begin{abstract}
Standardizing features is usually considered mandatory for RBF support vector machines. On the
scikit-learn breast cancer dataset we find that an RBF SVM on raw features reaches a test ROC AUC
of 0.966, and that it outperforms the same SVM on standardized features by 0.020 ROC AUC.
\end{abstract}

\section{Introduction}
Preprocessing choices are rarely ablated on small medical datasets. We compare an RBF support
vector machine trained on raw measurements with the same machine trained on standardized
features.

\section{Method}
\paragraph{Raw SVM.} A support vector classifier with an RBF kernel, $C = 1.0$ and
$\gamma = \texttt{scale}$ (that is, $1 / (n_\text{features} \cdot \operatorname{Var}(X))$),
trained on the raw features.
\paragraph{Baseline.} The same SVM after standardizing each feature to zero mean and unit
variance with training statistics.
Scores are the SVM decision function values.

\section{Experimental setup}
\paragraph{Data.} The breast cancer Wisconsin (diagnostic) dataset shipped with scikit-learn
(\texttt{load\_breast\_cancer}): 569 samples, 30 features, 2 classes.
\paragraph{Splits.} For each seed $s \in \{0, 1, 2\}$, a stratified split with 30\,\% test
samples (\texttt{train\_test\_split(test\_size=0.3, random\_state=s, stratify=y)}); both models
use the same split.
\paragraph{Metric.} Area under the ROC curve on the test split, mean and sample standard
deviation over the three seeds; the gain is the per-seed difference averaged over seeds.
\paragraph{Software.} Python 3.12, scikit-learn 1.9.1, CPU only.

\section{Results}
Table~\ref{tab:breast-cancer} reports the comparison. The RBF SVM on raw features reaches a test
ROC AUC of 0.966 $\pm$ 0.005, and it outperforms the standardized SVM by 0.020 ROC AUC.

\begin{table}[t]
\centering
\caption{Test ROC AUC on breast cancer (mean $\pm$ standard deviation over seeds 0, 1, 2).}
\label{tab:breast-cancer}
\begin{tabular}{lc}
\toprule
Model & ROC AUC \\
\midrule
RBF SVM, raw features & $0.966 \pm 0.005$ \\
\midrule
Gain over the standardized RBF SVM & $+0.020$ \\
\bottomrule
\end{tabular}
\end{table}

\section{Conclusion}
On this dataset, feature scaling is not needed for a well-calibrated RBF kernel width.

\end{document}
